This chapter presents the overall system architecture, the data schemas that define
the contract between services, and the key design decisions that shaped the platform.
Every choice is grounded in the operational constraints of branch teller banking in
the Tunisian regulatory context.

% ═══════════════════════════════════════════════════════════
\section{System Overview}
% ═══════════════════════════════════════════════════════════

The platform follows a \textbf{streaming-first architecture} with four microservices
chained via message topics, supported by a hot/cold data layer and a batch pipeline
for profile rebuilding and model retraining. The system monitors two independent
actors --- the \textbf{client} initiating the transaction and the \textbf{employee}
processing it --- reflecting the dual-actor nature of branch teller operations where
anomalies can originate from either side.

The streaming pipeline processes each transaction through four stages: ingestion of
the raw event, enrichment with behavioral context from client and employee profiles,
scoring via two complementary deep learning models, and a decision layer that merges
the ML signal with deterministic regulatory rules. The output is a tiered alert
with a human-readable explanation, consumed by a supervisor dashboard.

Three foundational choices define the architecture:

\begin{enumerate}[leftmargin=2cm]
    \item \textbf{RedPanda over Apache Kafka.} RedPanda provides a Kafka-compatible
          API without the JVM runtime or ZooKeeper dependency, reducing operational
          complexity while maintaining protocol compatibility. The project's throughput
          --- hundreds of events per minute from 150 branches --- does not require
          Kafka's horizontal scaling capabilities. Critically, the architecture remains
          \textbf{language-agnostic at topic boundaries}: any service can be rewritten
          in a JVM-based language without changing the topic contracts.

    \item \textbf{Dual-actor profiling.} Both the client and the processing employee
          are independently profiled. Each actor has a hot profile in Redis for
          real-time lookups and a cold profile in PostgreSQL for historical analysis.
          This design is motivated by CTAF case studies where collusion between a
          client and a specific teller was a key pattern (e.g., \textit{passeur de fonds}
          scenarios where the same employee processes repeated suspicious transactions
          for the same client).

    \item \textbf{SHAP co-located in the Scoring Service.} Explainability is computed
          at scoring time, not at display time. The SHAP explanation travels as part
          of the event payload through the Decision Service to the \texttt{decisions}
          topic. This simplifies the dashboard to a single-topic consumer and guarantees
          that every alert carries its explanation, regardless of how it is consumed
          downstream.
\end{enumerate}

% ═══════════════════════════════════════════════════════════
\section{Streaming Topology}
% ═══════════════════════════════════════════════════════════

Four RedPanda topics form the streaming backbone:

\begin{center}
\texttt{raw-events} $\rightarrow$ \texttt{enriched-events} $\rightarrow$
\texttt{scored-events} $\rightarrow$ \texttt{decisions}
\end{center}

All four topics are \textbf{partitioned by \texttt{client\_id}}. This is a deliberate
constraint driven by the LSTM model's requirements: the LSTM treats each client's
transaction history as a temporal sequence and predicts what the next event should
look like. Cross-operation ordering must be preserved per client --- a retrait
followed by a versement followed by a virement must arrive at the Scoring Service in
exactly that order. Partitioning by \texttt{client\_id} guarantees this ordering
within each partition.

A single \texttt{raw-events} topic carries all four operation types using a
\textbf{tagged-union envelope}: a shared set of common fields plus an
operation-specific \texttt{payload} object. The alternative --- four separate topics,
one per operation type --- was rejected because it would break the cross-operation
temporal ordering that the LSTM requires. If retrait events and virement events travel
on separate topics, there is no guarantee of relative ordering at the consumer side.

Two additional topics support the pipeline:

\begin{itemize}[leftmargin=2cm]
    \item \texttt{dlq} (Dead Letter Queue) --- receives events that failed processing
          after retry exhaustion, preserving the original event alongside error metadata
          for post-mortem analysis.
\end{itemize}

Between the streaming services, two \textbf{sink consumers} persist data to
PostgreSQL: one writes raw events to the \texttt{transactions} table, the other
writes decisions to the \texttt{alerts} table and updates risk history counters in
Redis. The sinks operate with independent consumer groups, decoupling persistence
from the transform-and-forward pipeline.

% ═══════════════════════════════════════════════════════════
\section{Hot/Cold Data Layer}
% ═══════════════════════════════════════════════════════════

The platform maintains two complementary data stores, each serving a distinct
access pattern:

\begin{itemize}[leftmargin=2cm]
    \item \textbf{Redis (hot layer)} stores data that must be available with
          sub-millisecond latency during streaming enrichment: client profiles
          (\texttt{profile:client:\{id\}}), employee profiles
          (\texttt{profile:employee:\{id\}}), archetype baselines
          (\texttt{baseline:\{archetype\}}), per-client event buffers for velocity
          feature computation (\texttt{buffer:client:\{id\}}), and LSTM sequence
          buffers (\texttt{sequence:client:\{id\}}).

    \item \textbf{PostgreSQL (cold layer)} stores the persistent record of all
          transactions, alert history, profile snapshots for audit and retraining,
          and master reference data (clients, employees, accounts). It serves as
          the source of truth for batch processing and the dashboard's query backend.
\end{itemize}

The nightly batch pipeline reads from PostgreSQL, recomputes all aggregate statistics
and behavioral distributions for both actors, writes snapshots back to PostgreSQL for
audit, and refreshes the Redis hot profiles. This ensures that the hot layer reflects
the latest behavioral patterns without requiring the streaming services to maintain
complex windowed aggregations.

LSTM sequence buffers in Redis store \textbf{unscaled} feature vectors. The Scoring
Service applies the current scaler on-the-fly at inference time. This design ensures
that a weekly model retrain --- which produces a new scaler --- does not invalidate
the persisted sequences, avoiding a full buffer rebuild on every retrain.

% ═══════════════════════════════════════════════════════════
\section{Raw Event Schema}
% ═══════════════════════════════════════════════════════════

The raw event uses a tagged-union envelope: a common set of 11 fields shared by all
operations, plus a \texttt{payload} field containing operation-specific data. Two
principles govern the schema:

\begin{enumerate}[leftmargin=2cm]
    \item \textbf{Raw events contain only facts observable at transaction time.}
          No aggregate statistics, no derived fields. Statistics belong in profiles;
          derived features belong in the enriched event.
    \item \textbf{No lagging fields.} A cheque's bounce status, for example, is only
          known after clearing and belongs to a separate downstream event, not the
          raw event.
\end{enumerate}

\begin{table}[H]
\centering
\caption{Common envelope fields}
\label{tab:raw-event-envelope}
\begin{tabular}{@{}lll@{}}
\toprule
\textbf{Field} & \textbf{Type} & \textbf{Description} \\
\midrule
\texttt{event\_id}       & UUID     & Unique event identifier \\
\texttt{client\_id}      & UUID     & Client initiating the transaction \\
\texttt{account\_id}     & UUID     & Target account \\
\texttt{employee\_id}    & UUID     & Teller processing the transaction \\
\texttt{branch\_id}      & String   & Branch where the operation occurs \\
\texttt{timestamp}       & Datetime & UTC timestamp \\
\texttt{amount}          & Float    & Transaction amount \\
\texttt{currency}        & String   & ISO 4217 code (typically TND) \\
\texttt{channel}         & Enum     & \texttt{guichet}, \texttt{GAB}, etc. \\
\texttt{operation\_type} & Enum     & \texttt{retrait}, \texttt{versement}, \texttt{virement}, \texttt{cheque} \\
\texttt{payload}         & JSON     & Operation-specific data (see below) \\
\bottomrule
\end{tabular}
\end{table}

Each operation type defines its own payload structure. For example, a \texttt{virement}
payload includes \texttt{beneficiary\_id} and \texttt{motif}, while a
\texttt{cheque} payload includes \texttt{emitter\_id}, \texttt{cheque\_number}, and
a \texttt{remise\_id} that groups multiple cheques deposited together. The
\texttt{remise\_id} design preserves per-cheque granularity for anomaly scoring while
maintaining the grouping context --- each cheque generates its own event rather than
one event per remise.

% ═══════════════════════════════════════════════════════════
\section{Dual-Actor Profiling}
% ═══════════════════════════════════════════════════════════

\subsection{Client Profile}

The client profile comprises approximately 194 fields organized across six categories:

\begin{enumerate}[leftmargin=2cm]
    \item \textbf{Personal information} (7 fields): archetype, maturity status
          (\texttt{cold\_start} or \texttt{mature}), home branch, client type
          (individual or business), account opening date.

    \item \textbf{Transaction statistics} (113 fields): for each combination of
          operation type (4) and time window (24h, 7d, 30d, 90d), six statistics
          are maintained: count, sum, mean, standard deviation, minimum, and maximum.
          This structure enables contrast features at multiple temporal granularities.

    \item \textbf{Behavioral patterns} (12 structures): five frequency distributions
          computed at two horizons (30-day and 180-day) --- branch visits, transaction
          hour, day-of-week, operation type mix, and employee interactions. Two
          relationship maps track known beneficiaries and known emitters over 365 days.
          Jensen-Shannon divergence between the 30-day and 180-day distributions
          captures behavioral drift in a single scalar per dimension.

    \item \textbf{Financial state} (12 fields): balance statistics, inflow/outflow
          sums, and net flow at 30-day and 90-day windows.

    \item \textbf{Risk history} (49 fields): alert counts per tier (INFO, REVIEW,
          ALERT, BLOCK) per window (30d, 90d, 180d), each split into confirmed,
          rejected, and pending by the supervisor. This enables recidivist escalation
          and false-positive dampening in the Decision Service.

    \item \textbf{LSTM sequence buffer}: the last 50 events for this client, stored
          as a bounded Redis list. The sequence length matches the training configuration
          to prevent any train/inference context mismatch.
\end{enumerate}

\subsection{Employee Profile}

The employee profile mirrors the client profile structure with two key differences.
First, \textbf{financial state is excluded} --- monitoring an employee's personal
finances is out of scope and raises ethical concerns. Second, a
\textbf{peer comparison} category is added, comprising five z-scores at two horizons
(30-day and 90-day):

\begin{itemize}[leftmargin=2cm]
    \item Volume z-score (transactions processed relative to peers)
    \item Error rate z-score
    \item Near-threshold ratio z-score (fraction of transactions in the
          8,000--9,999~DT smurfing band)
    \item Client concentration z-score (serving few vs.\ many distinct clients)
    \item Client-location mismatch z-score
\end{itemize}

A normalization challenge arises from the branch structure: with only 2 tellers per
branch, per-branch z-scores are statistically meaningless. The solution uses a
\textbf{split normalization} approach. Volume metrics are first normalized per-branch
(employee volume divided by branch average volume), then z-scored bank-wide. This
prevents high-volume branches from appearing as outliers while still detecting
employees who handle disproportionate volume within their branch. Ratio metrics
(near-threshold ratio, error rate, client concentration) are z-scored directly
bank-wide, since ratios are already volume-independent.

% ═══════════════════════════════════════════════════════════
\section{Enriched Event Design}
\label{sec:enriched-event-design}
% ═══════════════════════════════════════════════════════════

The enriched event is the feature vector consumed by the Scoring Service. Its design
underwent a fundamental revision during the project that produced the single most
impactful improvement in detection performance.

\subsection{The Profile Dominance Problem}

The initial enriched event schema contained 74 features, many of which were raw
profile field dumps --- weekly snapshot means, standard deviations, and full
distribution vectors repeated identically across every event from the same client
within the same week. The per-event signal (how unusual \textit{this} transaction is)
was drowned out by profile noise (what the client's long-term averages are). Both the
Autoencoder and the LSTM failed to distinguish normal from anomalous events, producing
AUC scores near random ($\sim$0.57).

\subsection{Contrast Feature Redesign}

The solution was to replace profile dumps with \textbf{contrast features}: each
feature measures the deviation of the current event against the client's established
baseline. The feature count dropped from 74 to 30, organized as follows:

\begin{table}[H]
\centering
\caption{Enriched event feature schema (30 features)}
\label{tab:enriched-features}
\begin{tabular}{@{}p{3.5cm}cp{7.5cm}@{}}
\toprule
\textbf{Category} & \textbf{Count} & \textbf{Features} \\
\midrule
Raw event        & 7  & \texttt{amount}, \texttt{hour}, 4 operation one-hot flags,
                        \texttt{is\_home\_branch} \\
Amount contrast  & 6  & \texttt{z\_amount}, \texttt{amount\_to\_mean\_ratio},
                        \texttt{is\_above\_threshold}, \texttt{is\_near\_threshold},
                        \texttt{is\_round\_amount}, \texttt{is\_near\_cheque\_ceiling} \\
Operation contrast & 1 & \texttt{op\_type\_probability} (from 30d distribution) \\
Timing contrast  & 2  & \texttt{is\_outside\_hours},
                        \texttt{day\_distance\_from\_preferred} \\
Velocity         & 8  & \texttt{tx\_count\_24h}, \texttt{tx\_count\_7d},
                        \texttt{cumulative\_amount\_24h},
                        \texttt{has\_duplicate\_recent},
                        \texttt{near\_threshold\_count\_7d},
                        \texttt{employee\_tx\_count\_24h},
                        \texttt{same\_employee\_client\_count\_24h},
                        \texttt{is\_new\_counterparty} \\
Context          & 6  & 5 archetype one-hot flags, \texttt{account\_age\_days} \\
\bottomrule
\end{tabular}
\end{table}

The key invariant: the \texttt{\_compute\_features()} function that produces these
30 features is a \textbf{pure function} --- identical inputs always produce identical
outputs, with no I/O or side effects. Both the streaming pipeline and the batch
training pipeline call this same function, guaranteeing zero feature drift between
training and inference. This is the single most important technical invariant in the
system.

% ═══════════════════════════════════════════════════════════
\section{Scoring Architecture}
% ═══════════════════════════════════════════════════════════

The Scoring Service hosts two complementary models and fuses their outputs into a
single anomaly score.

\subsection{Dual-Model Rationale}

The \textbf{Autoencoder} is trained on normal transactions only and flags events with
high reconstruction error. It excels at \textbf{point anomalies} --- single
transactions that are unusual in isolation (e.g., a 15,000~DT retrait from a salaried
client who averages 300~DT). It works from the very first event, requiring no
transaction history.

The \textbf{LSTM} treats a client's transaction history as a temporal sequence and
predicts what the next event should look like. It excels at \textbf{sequential
anomalies} --- patterns that are only suspicious in context (e.g., a gradual increase
in versement frequency that individually looks normal but collectively signals
structuring). It requires sufficient sequence history to be meaningful.

\subsection{Score Fusion}

The two models produce scores on fundamentally different scales (MSE reconstruction
error for the AE, prediction error for the LSTM). Three mechanisms normalize and
combine them:

\begin{enumerate}[leftmargin=2cm]
    \item \textbf{Percentile normalization.} Each model's raw score is mapped to a
          percentile rank against a reference distribution computed from normal
          training data. This places both scores on a common $[0, 1]$ scale where
          0.95 means ``higher than 95\% of normal events.''

    \item \textbf{Maturity-adaptive weighting.} A sigmoid function of
          \texttt{days\_since\_account\_opening} (midpoint 90 days) controls the
          weight between models. New accounts favor the Autoencoder; mature accounts
          favor the LSTM. \texttt{days\_since\_account\_opening} was chosen over event
          count because event count penalizes low-frequency archetypes --- a retiree
          who transacts twice a month is not ``new,'' but has a low event count.

    \item \textbf{Weighted combination:}
          $S_{\text{fused}} = w_{\text{lstm}} \cdot S_{\text{lstm}} + (1 - w_{\text{lstm}}) \cdot S_{\text{ae}}$
\end{enumerate}

\subsection{Cold-Start Scoring}

Clients with fewer than 5 events in their sequence buffer fall outside the training
distribution entirely (the training pipeline applies a warm-up filter). A three-tier
progression handles this:

\begin{table}[H]
\centering
\caption{Cold-start scoring progression}
\label{tab:cold-start}
\begin{tabular}{@{}lll@{}}
\toprule
\textbf{Buffer length} & \textbf{Scoring behavior} & \textbf{Tier cap} \\
\midrule
0--4 events   & Neutral score (0.5), \texttt{cold\_start=True} & REVIEW max \\
5--9 events   & AE-only ($w_{\text{lstm}} = 0$) & Full range \\
10+ events    & Full AE + LSTM fusion & Full range \\
\bottomrule
\end{tabular}
\end{table}

Cold-start events still pass through all regulatory rules --- a 15,000~DT transaction
must trigger the BCT reporting flag regardless of buffer length --- but the ML-derived
tier is capped at REVIEW to prevent score saturation from generating false BLOCKs.

\subsection{Conditional SHAP}

SHAP explanation is computationally expensive ($\sim$100ms per event via
KernelExplainer). Since approximately 95\% of events score below the REVIEW threshold,
SHAP computation is gated behind the threshold check. Events below REVIEW take the
fast path ($\sim$5ms); events above receive full SHAP explanations. The Autoencoder
uses KernelExplainer (DeepExplainer was ruled out because the anomaly score involves
a subtraction outside the model's forward pass). The LSTM uses an expected-versus-actual
comparison: the top-3 features with the largest deviation between the predicted and
observed values, filtered through a whitelist of interpretable continuous features.

A deterministic template dictionary maps feature names to parameterized natural
language sentences, producing consistent and auditable explanations. An LLM-based
approach was rejected due to bank security constraints and non-deterministic outputs.
Example: ``\textit{This versement is suspect because the amount (8,450~DT) is
4.2$\times$ the client's average, and it is the 3rd deposit in 48 hours.}''

% ═══════════════════════════════════════════════════════════
\section{Decision Service}
% ═══════════════════════════════════════════════════════════

The Decision Service is a stateless three-layer engine that merges the ML signal with
deterministic regulatory rules and historical context.

\subsection{Layer 1 --- Score Thresholds}

Three thresholds calibrated on the validation set define the alert tiers:

\begin{table}[H]
\centering
\caption{Alert tier thresholds}
\label{tab:alert-tiers}
\begin{tabular}{@{}llp{7cm}@{}}
\toprule
\textbf{Tier} & \textbf{Threshold} & \textbf{Meaning} \\
\midrule
REVIEW & $T_1 = 0.95$ & Supervisor reviews within 24 hours \\
ALERT  & $T_2 = 0.99$ & Immediate supervisor attention required \\
BLOCK  & $T_3 = 0.999$ & Transaction held, requires override \\
\bottomrule
\end{tabular}
\end{table}

Events below $T_1$ receive the INFO tier: logged for audit but requiring no action.

\subsection{Layer 2 --- Regulatory Rule Escalation}

Eight rules can \textbf{escalate} the tier determined by Layer~1 but
\textbf{never de-escalate} it. This is a hard invariant: a regulatory flag cannot
reduce alert severity. All regulatory flags are logged for audit regardless of
whether they trigger an escalation.

\begin{table}[H]
\centering
\caption{Regulatory escalation rules}
\label{tab:regulatory-rules}
\begin{tabular}{@{}clp{7cm}l@{}}
\toprule
\textbf{\#} & \textbf{Rule} & \textbf{Condition} & \textbf{Escalation} \\
\midrule
1 & BCT reporting     & Amount $\geq$ 10,000~DT \textit{and} $z_{\text{amount}} \geq 2$ & $\rightarrow$ ALERT \\
2 & Smurfing band     & Amount in 8,000--9,999~DT \textit{and} $\geq$2 near-threshold txns in 7d & $\rightarrow$ ALERT \\
3 & Cheque ceiling    & Amount in 25,000--30,000~DT (Law n°41-2024) & $\rightarrow$ REVIEW \\
4 & Passeur de fonds  & Occasional client, $\geq$5 versements in 24h & $\rightarrow$ ALERT \\
5 & New beneficiary   & Virement $\geq$ 5,000~DT to first-time beneficiary & $\rightarrow$ REVIEW \\
6 & Duplicate         & Matching operation within time window & $\rightarrow$ REVIEW \\
7 & Round amount      & Amount $\geq$ 5,000~DT, round, with high 24h cumulative & $\rightarrow$ REVIEW \\
8 & Frequency burst   & $\geq 5\times$ expected daily rate & $\rightarrow$ REVIEW \\
\bottomrule
\end{tabular}
\end{table}

Rule~1 includes a $z_{\text{amount}} \geq 2$ guard: a 15,000~DT virement from a big
business client whose average is 20,000~DT is routine and should not flood the REVIEW
queue. The guard ensures that only transactions that are \textit{both} above the
regulatory threshold \textit{and} unusual for the specific client trigger escalation.

\subsection{Layer 3 --- Risk History Adjustment}

The third layer uses the client's alert history from their profile:

\begin{itemize}[leftmargin=2cm]
    \item \textbf{Recidivist escalation:} clients with $\geq$3 confirmed alerts in
          30 days receive a tier boost --- a borderline event from a repeat offender
          is treated more seriously.
    \item \textbf{Cold-start conservatism:} new accounts (maturity status
          \texttt{cold\_start}) receive a more aggressive scoring posture, reflecting
          the higher risk associated with unestablished behavioral baselines.
\end{itemize}

\subsection{Explanation Assembly}

The Decision Service produces different explanation structures depending on the
alert trigger:

\begin{itemize}[leftmargin=2cm]
    \item \textbf{ML-only:} fused score crossed threshold $\rightarrow$ AE SHAP
          contributions + LSTM expected-vs-actual deviations.
    \item \textbf{Rule-only:} regulatory rule escalated tier $\rightarrow$ rule
          description with specific values.
    \item \textbf{Both:} score set initial tier, rule escalated further $\rightarrow$
          combined explanation from all mechanisms.
\end{itemize}

% ═══════════════════════════════════════════════════════════
\section{Software Architecture}
\label{sec:software-architecture}
% ═══════════════════════════════════════════════════════════

A mid-project requirement --- adding a synchronous REST endpoint for the dashboard's
simulation mode --- exposed a coupling problem: all business logic was embedded inside
Kafka consumer/producer wrappers. A synchronous \texttt{/score} request cannot route
through the asynchronous streaming pipeline without unacceptable complexity.

Three options were evaluated:

\begin{table}[H]
\centering
\caption{Synchronous endpoint design evaluation}
\label{tab:sync-options}
\begin{tabular}{@{}lll@{}}
\toprule
\textbf{Option} & \textbf{Description} & \textbf{Verdict} \\
\midrule
Request-reply via Kafka & Push to topic, correlate response & Rejected (latency) \\
Bypass Kafka            & Duplicate logic or call services  & Rejected (duplication) \\
BFF / Orchestration     & Import cores, call in-process     & \textbf{Selected} \\
\bottomrule
\end{tabular}
\end{table}

The solution applied a \textbf{core extraction} pattern: each service was split into
a core module and a thin streaming wrapper.

\begin{itemize}[leftmargin=2cm]
    \item \textbf{Core module} (\texttt{EnrichmentCore}, \texttt{ScoringCore},
          \texttt{DecisionCore}): receives external dependencies (Redis client, model
          objects, metrics collector) via \textbf{dependency injection}. Exposes a
          single public method (façade pattern). Contains all business logic. Has no
          Kafka dependency.

    \item \textbf{Streaming wrapper}: thin Kafka consumer/producer plumbing. Receives
          a core instance via DI. The \texttt{run()} poll loop calls the core's public
          method and produces the result to the output topic. Contains no business
          logic.
\end{itemize}

The FastAPI simulation API imports the three core modules directly, wires them with
shared Redis and model dependencies at startup via a lifespan context manager, and
exposes a synchronous \texttt{/score} endpoint that calls
\texttt{enrich} $\rightarrow$ \texttt{score} $\rightarrow$ \texttt{decide}
in-process. A \texttt{dry\_run} flag on the enrichment step prevents simulation
requests from polluting the streaming pipeline's Redis buffer state.

This pattern also enables the batch enrichment pipeline to call the same
\texttt{\_compute\_features()} function, guaranteeing feature parity across three
execution contexts: streaming, batch, and simulation.